\chapter{การออกแบบระบบ แหล่งจ่ายไฟฟ้าพัลส์ และความปลอดภัย}
\label{chap:ch10}

% =========================================================================
% แผนบริหารการสอนประจำบทที่ 10 (Modular Lesson Plan Boxes - 2-Page Standard)
% =========================================================================

\vspace*{-6pt}
\section*{แผนบริหารการสอนประจำบทที่ 10}
\addcontentsline{toc}{section}{แผนบริหารการสอนประจำบทที่ 10}

% --- Page 1 of Lesson Plan ---
\begin{planbox}{หัวข้อเนื้อหาประจำบท}
\begin{enumerate}
    \item สถาปัตยกรรมแหล่งจ่ายไฟฟ้าแรงสูงสำหรับพลาสมาเย็นความดันบรรยากาศ
    \item โทโพโลยีวงจรพัลส์กำลังสูง: มาร์กซ์เจเนอเรเตอร์สารกึ่งตัวนำ (Solid-State Marx Generator)
    \item การแมตชิ่งอิมพีแดนซ์พลวัต (Dynamic Impedance Matching) ของโหลดพลาสมา
    \item สวิตช์สารกึ่งตัวนำแถบพลังงานกว้าง (Wide Bandgap Semiconductors: SiC และ GaN)
    \item การลดทอนสัญญาณรบกวนทางแม่เหล็กไฟฟ้า (Electromagnetic Interference หรือ EMI)
    \item สถาปัตยกรรม Edge AI Digital Twin สำหรับการพยากรณ์ความผิดพร่องและการป้องกันอาร์ก
    \item มาตรฐานความปลอดภัยทางวิศวกรรมไฟฟ้าแรงสูงและการจัดการก๊าซโอโซนในแปลงเกษตร
\end{enumerate}
\end{planbox}

\begin{planbox}{วัตถุประสงค์เชิงพฤติกรรม}
\begin{enumerate}
    \item วิเคราะห์โทโพโลยีวงจรเรียงกระแส อินเวอร์เตอร์เรโซแนนซ์ และมาร์กซ์เจเนอเรเตอร์ได้
    \item คำนวณพลังงานสะสมต่อพัลส์ กำลังไฟฟ้าเฉลี่ย และอัตราการเพิ่มของแรงดัน ($\text{d}V/\text{d}t$) ได้
    \item ออกแบบวงจรกรองและสายดินเพื่อลดสัญญาณรบกวน EMI ตามมาตรฐาน IEC 61000 ได้
    \item อธิบายการทำงานของโครงข่ายประสาทเทียมตรวจจับความผิดปกติแบบออโตเอ็นโคเดอร์ (Autoencoder) ได้
    \item ประเมินความเสี่ยงและกำหนดมาตรการความปลอดภัยทางไฟฟ้าและก๊าซพิษในแปลงเกษตรได้
\end{enumerate}
\end{planbox}

\newpage

% --- Page 2 of Lesson Plan ---
\begin{planbox}{วิธีสอนและกิจกรรมการเรียนการสอน}
\begin{enumerate}
    \item บรรยายเชิงทฤษฎีเกี่ยวกับวิศวกรรมไฟฟ้าแรงสูงและอิเล็กทรอนิกส์กำลังสำหรับพลาสมา
    \item สาธิตการทำงานของวงจรอินเวอร์เตอร์เรโซแนนซ์แบบฟูลบริดจ์และการจำลองด้วยโปรแกรม SPICE
    \item ปฏิบัติการวัดรูปคลื่นแรงดันสูงด้วย High-Voltage Probe และโพรบวัดกระแส Rogowski Coil
    \item ปฏิบัติการทดสอบค่าความเป็นฉนวนไฟฟ้าและความต้านทานสายดินตามมาตรฐานสากล
    \item ฝึกทดสอบโมเดล Edge AI Autoencoder ในการตรวจจับความผิดปกติของรูปคลื่นกระแสไฟฟ้า
\end{enumerate}
\end{planbox}

\begin{planbox}{สื่อการเรียนการสอน}
\begin{enumerate}
    \item แหล่งจ่ายไฟฟ้าพัลส์นาโนวินาทีแรงดันสูง ($0 - 20 \text{ kV}$, ความถี่ $1 - 100 \text{ kHz}$)
    \item ออสซิลโลสโคปความเร็วสูง (แบนด์วิดท์ $\ge 500 \text{ MHz}$) พร้อมโพรบแรงดันสูง $1000:1$
    \item เครื่องมือวัดความเข้ากันได้ทางแม่เหล็กไฟฟ้า (EMI Receiver / Spectrum Analyzer)
    \item บอร์ดประมวลผลสมองกลฝังตัวประสิทธิภาพสูงสำหรับการทำ Digital Twin
    \item คู่มือมาตรฐานความปลอดภัยทางวิศวกรรมไฟฟ้าแรงสูงและการระบายอากาศ
\end{enumerate}
\end{planbox}

\begin{planbox}{การวัดและประเมินผล}
\begin{enumerate}
    \item ประเมินทักษะการออกแบบและต่อวงจรไฟฟ้าแรงสูงอย่างปลอดภัย ร้อยละ 20
    \item ประเมินรายงานผลการวัดรูปคลื่นพัลส์และการคำนวณกำลังงานสูญเสีย ร้อยละ 20
    \item ประเมินการทำแบบฝึกหัดการคำนวณวงจรเรโซแนนซ์และการตัดทอน EMI ร้อยละ 10
    \item ประเมินผลการสอบวัดความรู้เชิงทฤษฎีประจำภาคการศึกษา ร้อยละ 50
\end{enumerate}
\end{planbox}

\clearpage

% =========================================================================
% เนื้อหาบทที่ 10
% =========================================================================

\section{วิศวกรรมแหล่งจ่ายไฟฟ้าแรงสูงสำหรับพลาสมาความดันบรรยากาศ}

หัวใจสำคัญในการขับเคลื่อนระบบพลาสมาเย็นความดันบรรยากาศคือ แหล่งจ่ายไฟฟ้าแรงสูง (High-Voltage Power Supplies) โหลดพลาสมา โดยเฉพาะในการคายประจุแบบแผงกั้นไดอิเล็กทริก (DBD) หรือพลาสมาเจ็ท (APPJ) มีลักษณะทางไฟฟ้าที่เป็นโหลดแบบไม่เป็นเชิงเส้นและแปรผันตามเวลา (Non-linear Time-Varying Load) ก่อนการเบรกดาวน์ โหลดจะประพฤติตัวเสมือนเป็นตัวเก็บประจุบริสุทธิ์ของไดอิเล็กทริกและช่องว่างก๊าซ แต่ทันทีที่เกิดการแตกตัวเป็นไอออน ช่องว่างก๊าซจะกลายเป็นตัวต้านทานที่มีค่าความนำไฟฟ้าสูงอย่างฉับพลันภายในเวลาไม่กี่นาโนวินาที

แหล่งจ่ายไฟฟ้าที่ใช้ในงานเกษตรกรรมและสิ่งแวดล้อมแบ่งออกเป็น 3 สถาปัตยกรรมหลัก ได้แก่

\begin{enumerate}
    \item \textbf{แหล่งจ่ายไฟฟ้ากระแสสลับไซนูซอยด์ความถี่สูง (High-Frequency AC Resonant):} ใช้ความถี่ในช่วง $10 - 100 \text{ kHz}$ แรงดัน $5 - 25 \text{ kV}_{\text{p-p}}$ โดยอาศัยหลักการเรโซแนนซ์แบบอนุกรมหรือขนานระหว่างตัวเก็บประจุของโหลด DBD กับตัวเหนี่ยวนำรั่วไหลของหม้อแปลงแรงดันสูง (Leakage Inductance) เพื่อลดการสูญเสียจากการสวิตชิ่ง (Zero Voltage Switching หรือ ZVS)
    
    \item \textbf{แหล่งจ่ายไฟฟ้าพัลส์นาโนวินาทีแรงดันสูง (Nanosecond High-Voltage Pulsed Power):} ป้อนพัลส์แรงดันสูงที่มีระยะเวลาการเพิ่มของแรงดัน (Rise Time) สั้นกว่า $10 - 50 \text{ ns}$ และความกว้างพัลส์ $50 - 200 \text{ ns}$ การใช้พัลส์นาโนวินาทีมีข้อได้เปรียบทางอุณหพลศาสตร์อย่างมหาศาล เนื่องจากสนามไฟฟ้าที่สูงจัดจะส่งพลังงานให้แก่อิเล็กตรอนโดยตรง ก่อนที่ไอออนหนักจะทันเคลื่อนที่ ส่งผลให้ก๊าซยังคงมีอุณหภูมิห้อง ป้องกันการเกิดความร้อนสะสมที่อาจทำลายเนื้อเยื่อพืชหรือเมล็ดพันธุ์
    
    \item \textbf{แหล่งจ่ายไฟฟ้าพัลส์ความถี่วิทยุ (RF Power Supplies):} ทำงานที่ความถี่มาตรฐานอุตสาหกรรม $13.56 \text{ MHz}$ ให้พลาสมาที่มีความหนาแน่นอิเล็กตรอนสูงและความสม่ำเสมอเป็นเลิศ แต่วงจรแมตชิ่งมีความซับซ้อนและมีต้นทุนสูง
\end{enumerate}

\section{โทโพโลยีวงจรมาร์กซ์เจเนอเรเตอร์สารกึ่งตัวนำ}

ในอดีต วงจรมาร์กซ์เจเนอเรเตอร์ดั้งเดิมใช้ช่องว่างประกายไฟ (Spark Gaps) เป็นสวิตช์ ซึ่งมีข้อจำกัดด้านความถี่การทำซ้ำต่ำ เกิดการสึกหรอของขั้วไฟฟ้า และมีสัญญาณรบกวนแม่เหล็กไฟฟ้ารุนแรง ในปัจจุบัน การพัฒนาสวิตช์สารกึ่งตัวนำแถบพลังงานกว้าง (Wide Bandgap Semiconductors) เช่น ซิลิคอนคาร์ไบด์ (Silicon Carbide หรือ SiC MOSFETs) และแกลเลียมไนไตรด์ (Gallium Nitride หรือ GaN HEMTs) ช่วยให้สามารถสร้าง ``มาร์กซ์เจเนอเรเตอร์สารกึ่งตัวนำ (Solid-State Marx Generator หรือ SSMG)'' ที่ทำงานด้วยความถี่สูงหลายสิบกิโลเฮิรตซ์ได้อย่างมีเสถียรภาพ

\subsection{หลักการประจุแบบขนานและคายประจุแบบอนุกรม}
วงจร SSMG ประกอบด้วยสเตจย่อย $N$ สเตจ แต่ละสเตจประกอบด้วยตัวเก็บประจุสะสมพลังงาน ($C_{\text{stage}}$), สวิตช์หลัก (SiC MOSFET, $S_{\text{main}}$) และสวิตช์ประจุหรือไดโอดความเร็วสูง ($D_{\text{charge}}$)

\begin{itemize}
    \item \textbf{สถานะการประจุ (Charging State):} ตัวเก็บประจุทุกตัวในแต่ละสเตจจะถูกต่อขนานกันเพื่อรับแรงดันไฟฟ้ากระแสตรงจากแหล่งจ่าย $V_{\text{DC}}$ (เช่น $1000 \text{ V}$) ผ่านสวิตช์ประจุ
    \item \textbf{สถานะการคายประจุ (Discharging State):} สวิตช์หลัก $S_{\text{main}}$ ทุกตัวจะถูกสั่งให้เปิดวงจรพร้อมกันผ่านตัวขับเกตแบบแยกกราวด์ (Isolated Gate Drivers) ส่งผลให้ตัวเก็บประจุทุกตัวถูกต่อเรียงกันแบบอนุกรม เกิดเป็นพัลส์แรงดันสูงขาออกรวม
    \begin{equation}
        V_{\text{pulse}} = N \cdot V_{\text{DC}}
        \label{eq:marx_voltage}
    \end{equation}
\end{itemize}

พลังงานที่สะสมต่อพัลส์ ($E_{\text{pulse}}$) และกำลังไฟฟ้าเฉลี่ย ($P_{\text{avg}}$) คำนวณได้จากสมการ

\begin{align}
    E_{\text{pulse}} &= \frac{1}{2} C_{\text{total}} V_{\text{pulse}}^2 = \frac{1}{2} \left( \frac{C_{\text{stage}}}{N} \right) (N \cdot V_{\text{DC}})^2 = \frac{1}{2} N C_{\text{stage}} V_{\text{DC}}^2 \label{eq:marx_energy} \\
    P_{\text{avg}} &= E_{\text{pulse}} \cdot f_{\text{rep}} = \frac{1}{2} N C_{\text{stage}} V_{\text{DC}}^2 f_{\text{rep}} \label{eq:marx_power}
\end{align}

โดยที่ $f_{\text{rep}}$ คือความถี่การเกิดพัลส์ซ้ำ (Pulse Repetition Frequency, Hz)

\section{การแมตชิ่งอิมพีแดนซ์พลวัตและการถ่ายโอนพลังงาน}

ประสิทธิภาพการถ่ายโอนพลังงานจากแหล่งจ่ายไฟฟ้าพัลส์ไปยังโหลดพลาสมาขึ้นอยู่กับความสอดคล้องของอิมพีแดนซ์ (Impedance Matching) ค่าอิมพีแดนซ์ลักษณะเฉพาะของสายส่งและแหล่งจ่าย ($Z_0$) ต้องมีค่าใกล้เคียงกับอิมพีแดนซ์ของโหลดพลาสมา ($Z_{\text{load}}$)

สัมประสิทธิ์การสะท้อนกลับของแรงดัน ($\Gamma$) และกำลังงานที่สะท้อนกลับ ($P_{\text{reflected}}$) อธิบายด้วยสมการ

\begin{equation}
    \Gamma = \frac{Z_{\text{load}} - Z_0}{Z_{\text{load}} + Z_0} \quad \implies \quad \frac{P_{\text{reflected}}}{P_{\text{incident}}} = |\Gamma|^2
    \label{eq:reflection_coeff}
\end{equation}

เนื่องจาก $Z_{\text{load}}$ เปลี่ยนแปลงอย่างรุนแรงระหว่างช่วงก่อนเบรกดาวน์ ($|Z_{\text{load}}| \to \infty$) และช่วงเบรกดาวน์ ($|Z_{\text{load}}| \approx 50 - 200 \ \Omega$) การออกแบบวงจรจึงต้องมีการใส่เครือข่ายชดเชยเรโซแนนซ์แบบแอลซี ($L\text{-}C$ Matching Network) หรือการใช้วงจรลดทอนแรงดันย้อนกลับ (Snubber Circuits) เพื่อป้องกันแรงดันกระชาก (Voltage Spikes) ไม่ให้ย้อนกลับมาทำลายสวิตช์สารกึ่งตัวนำ

\section{การควบคุมสัญญาณรบกวนแม่เหล็กไฟฟ้าและการกราวดิ้ง}

การคายประจุพลาสมาความดันบรรยากาศสร้างการเปลี่ยนแปลงของกระแสและแรงดันอย่างฉับพลัน ($\text{d}i/\text{d}t > 10^9 \text{ A/s}$ และ $\text{d}v/\text{d}t > 10^{11} \text{ V/s}$) ซึ่งเป็นแหล่งกำเนิดสัญญาณรบกวนทางแม่เหล็กไฟฟ้า (Electromagnetic Interference หรือ EMI) ย่านกว้าง ตั้งแต่ความถี่วิทยุไม่กี่กิโลเฮิรตซ์จนถึงย่านกิกะเฮิรตซ์ หากไม่มีการออกแบบการชีลด์และการกราวดิ้งที่ถูกต้อง สัญญาณรบกวนจะแพร่กระจายไปรบกวนระบบสมองกลฝังตัว เซนเซอร์ในแปลงเกษตร หรือระบบสื่อสารไร้สาย (Wi-Fi/4G)

หลักการออกแบบเพื่อความเข้ากันได้ทางแม่เหล็กไฟฟ้า (Electromagnetic Compatibility หรือ EMC) ตามมาตรฐาน IEC 61000 ประกอบด้วย

\begin{enumerate}
    \item \textbf{การออกแบบตู้ป้องกันแบบกรงฟาราเดย์ (Faraday Cage Enclosure):} ตัวถังตู้ควบคุมพลาสมาต้องทำจากโลหะนำไฟฟ้า (อะลูมิเนียมหรือสแตนเลสหนา $\ge 1.5 \text{ mm}$) รอยต่อทุกจุดต้องมีปะเก็นตัวนำไฟฟ้า (Conductive Gaskets) เพื่อป้องกันการรั่วไหลของสนามแม่เหล็กไฟฟ้า
    \item \textbf{การแยกกราวด์อย่างเด็ดขาด (Ground Isolation):} แบ่งระบบกราวด์ออกเป็น 3 ระบบที่แยกจากกัน ได้แก่ กราวด์ไฟฟ้าแรงสูง (HV Ground), กราวด์วงจรดิจิทัล/ไมโครคอนโทรลเลอร์ (Digital Signal Ground), และกราวด์เพื่อความปลอดภัย (Chassis Earth Ground) โดยเชื่อมต่อกันที่จุดเดียว (Single-Point Star Grounding)
    \item \textbf{การใช้ตัวแยกสัญญาณเชิงแสง (Optocouplers / Fiber Optics):} สัญญาณสั่งการทริกเกอร์เกตระหว่างบอร์ดควบคุมดิจิทัลกับสวิตช์แรงดันสูงต้องส่งผ่านสายใยแก้วนำแสง (Optical Fibers) เท่านั้น เพื่อตัดวงจรลูปกราวด์โดยสิ้นเชิง
\end{enumerate}

\section{ความปลอดภัยทางวิศวกรรมไฟฟ้าและสารเคมีในแปลงเกษตรกรรม}

การนำระบบพลาสมาไปใช้งานจริงในแปลงเกษตรกรรมหรือโรงคัดบรรจุต้องคำนึงถึงสภาพแวดล้อมที่เปียกชื้น มีฝุ่นละออง และผู้ใช้งานเป็นเกษตรกรทั่วไป มาตรการความปลอดภัยต้องเป็นไปตามมาตรฐานวิศวกรรมไฟฟ้าสากล

\begin{itemize}
    \item \textbf{มาตรฐานการป้องกันน้ำและฝุ่น (Ingress Protection):} ตู้ควบคุมภาคสนามต้องผ่านมาตรฐานระดับ $\text{IP65}$ หรือ $\text{IP66}$ ขึ้นไป
    \item \textbf{ระบบตัดไฟฉุกเฉินและอินเตอร์ล็อก (Fail-Safe Interlock System):} ประตูตู้ปฏิกรณ์พลาสมาต้องมีสวิตช์อินเตอร์ล็อกแม่เหล็ก เมื่อเปิดประตู วงจรไฟฟ้าแรงสูงจะถูกตัดไฟทันทีภายใน $10 \ \mu\text{s}$ และมีรีเลย์คายประจุตกค้างจากตัวเก็บประจุลงสู่สายดินโดยอัตโนมัติ
    \item \textbf{ระบบตรวจจับและทำลายก๊าซโอโซนตกค้าง (Ozone Safety System):} โอโซนเป็นก๊าซพิษต่อระบบทางเดินหายใจ โดยมาตรฐาน OSHA กำหนดขีดจำกัดความเข้มข้นในสถานที่ทำงานไม่เกิน $0.10 \text{ ppm}$ ในระบบพลาสมาจึงต้องติดตั้งตัวทำลายโอโซนแบบเร่งปฏิกิริยา ($\ce{MnO2/CuO}$ Catalyst Destructor) เพื่อสลายโอโซนส่วนเกินให้กลายเป็นออกซิเจน ($\ce{O2}$) ก่อนปล่อยสู่บรรยากาศภายนอก
\end{itemize}

% -------------------------------------------------------------------------
% Worked Calculation Example Box
% -------------------------------------------------------------------------
\begin{examplebox}{การคำนวณพารามิเตอร์วงจรเรโซแนนซ์และกำลังงานสูญเสียในสวิตช์ SiC}
\textbf{โจทย์คำนวณ:} แหล่งจ่ายไฟฟ้าความถี่สูงแบบเรโซแนนซ์ชุดหนึ่งถูกออกแบบมาเพื่อขับเคลื่อนเครื่องปฏิกรณ์ DBD สำหรับกระตุ้นเมล็ดพันธุ์ข้าว โหลด DBD มีความจุไฟฟ้าของไดอิเล็กทริก $C_{\text{d}} = 800 \text{ pF}$ และความจุไฟฟ้าของช่องว่างก๊าซ $C_{\text{g}} = 200 \text{ pF}$ วงจรอินเวอร์เตอร์แบบฟูลบริดจ์ใช้สวิตช์ SiC MOSFET ทำงานที่ความถี่เรโซแนนซ์ $f_0 = 40.0 \text{ kHz}$ โดยต่อตัวเหนี่ยวนำภายนอก ($L_{\text{ext}}$) อนุกรมกับหม้อแปลงแรงดันสูงที่มีอัตราทดรอบ $N_1:N_2 = 1:20$ และความเหนี่ยวนำรั่วไหลฝั่งทุติยภูมิ $L_{\text{leak}} = 15.0 \text{ mH}$ แรงดันป้อนเข้าฝั่งปฐมภูมิ $V_{\text{DC}} = 300 \text{ V}$ จงคำนวณ
\begin{enumerate}
    \item ความจุไฟฟ้ารวมของโหลด DBD ในช่วงก่อนเบรกดาวน์ ($C_{\text{cell}}$)
    \item ความเหนี่ยวนำรวมที่ต้องการ ($L_{\text{total}}$) ในฝั่งทุติยภูมิเพื่อให้ระบบเกิดเรโซแนนซ์ที่ $f_0 = 40.0 \text{ kHz}$
    \item ค่าความเหนี่ยวนำภายนอกที่ต้องติดตั้งเพิ่ม ($L_{\text{ext}}$) ในฝั่งทุติยภูมิ
    \item หากแรงดันยอดคลื่นฝั่งทุติยภูมิมีค่า $V_{\text{peak}} = 12.0 \text{ kV}$ จงคำนวณพลังงานสะสมสูงสุดในตัวเก็บประจุ ($E_{\text{cap}}$) และกำลังไฟฟ้าเฉลี่ยเมื่อป้อนความถี่ 40 kHz
    \item คำนวณกำลังงานสูญเสียจากการนำกระแส (Conduction Loss) ในสวิตช์ SiC 4 ตัวในบริดจ์ กำหนดให้ความต้านทานขณะเปิด $R_{\text{DS(on)}} = 45 \text{ m}\Omega$ และกระแสประสิทธิผลฝั่งปฐมภูมิ $I_{\text{rms}} = 12.5 \text{ A}$
\end{enumerate}

\vspace{6pt}
\textbf{วิธีทำ:}

\textbf{ขั้นที่ 1: คำนวณความจุไฟฟ้ารวมของโหลด DBD ก่อนเบรกดาวน์}
เนื่องจาก $C_{\text{d}}$ และ $C_{\text{g}}$ ต่ออนุกรมกันทางกายภาพ:
\begin{equation*}
    C_{\text{cell}} = \frac{C_{\text{d}} \cdot C_{\text{g}}}{C_{\text{d}} + C_{\text{g}}} = \frac{800 \times 200}{800 + 200} = \frac{160,000}{1,000} = 160 \text{ pF} = 160 \times 10^{-12} \text{ F}
\end{equation*}

\textbf{ขั้นที่ 2: คำนวณความเหนี่ยวนำรวมที่ต้องการที่ความถี่เรโซแนนซ์}
สมการความถี่เรโซแนนซ์ของทอมสัน:
\begin{align*}
    \omega_0 &= 2 \pi f_0 = 2 \times 3.1416 \times 40,000 = 251,327 \text{ rad/s} \\
    L_{\text{total}} &= \frac{1}{\omega_0^2 C_{\text{cell}}} = \frac{1}{(251,327)^2 \times (160 \times 10^{-12})} \\
    &= \frac{1}{(6.3165 \times 10^{10}) \times (1.60 \times 10^{-10})} = \frac{1}{10.106} = 0.09895 \text{ H} \approx 98.95 \text{ mH}
\end{align*}

\textbf{ขั้นที่ 3: คำนวณค่าความเหนี่ยวนำภายนอกที่ต้องติดตั้งเพิ่ม}
\begin{equation*}
    L_{\text{ext}} = L_{\text{total}} - L_{\text{leak}} = 98.95 \text{ mH} - 15.00 \text{ mH} = 83.95 \text{ mH}
\end{equation*}

\textbf{ขั้นที่ 4: คำนวณพลังงานสะสมและกำลังไฟฟ้า}
\begin{align*}
    E_{\text{cap}} &= \frac{1}{2} C_{\text{cell}} V_{\text{peak}}^2 = \frac{1}{2} \times (160 \times 10^{-12} \text{ F}) \times (12,000 \text{ V})^2 \\
    &= 80 \times 10^{-12} \times 1.44 \times 10^8 = 0.01152 \text{ J} = 11.52 \text{ mJ/cycle}
\end{align*}
กำลังไฟฟ้าสั่นพ้องสูงสุดตามทฤษฎี:
\begin{equation*}
    P_{\text{theo}} = E_{\text{cap}} \times f_0 = 0.01152 \text{ J} \times 40,000 \text{ s}^{-1} = 460.8 \text{ W}
\end{equation*}

\textbf{ขั้นที่ 5: คำนวณกำลังงานสูญเสียในสวิตช์ SiC MOSFETs}
สำหรับอินเวอร์เตอร์แบบฟูลบริดจ์ (4 สวิตช์):
\begin{align*}
    P_{\text{cond, switch}} &= I_{\text{rms}}^2 \times R_{\text{DS(on)}} \times 0.5 \quad (\text{แต่ละตัวนำกระแสครึ่งคาบ}) \\
    P_{\text{cond, total}} &= 4 \times \left( I_{\text{rms}}^2 \times R_{\text{DS(on)}} \times 0.5 \right) = 2 \times I_{\text{rms}}^2 \times R_{\text{DS(on)}} \\
    &= 2 \times (12.5 \text{ A})^2 \times (0.045 \ \Omega) = 2 \times 156.25 \times 0.045 = 14.06 \text{ W}
\end{align*}

\textbf{สรุปผลการวิเคราะห์:} การเลือกใช้สวิตช์ SiC MOSFET ร่วมกับวงจรเรโซแนนซ์ที่คำนวณได้ ($L_{\text{ext}} = 83.95 \text{ mH}$) ส่งผลให้กำลังงานสูญเสียในการนำกระแสของสวิตช์มีค่าต่ำเพียง $14.06 \text{ W}$ เมื่อเทียบกับกำลังขับเคลื่อนรวม $460 \text{ W}$ (คิดเป็นสัดส่วนการสูญเสียเพียงร้อยละ 3.05) ทำให้ระบบมีประสิทธิภาพทางพลังงานสูงกว่าร้อยละ 95 และมีความร้อนสะสมในตู้คอนโทรลต่ำมาก
\end{examplebox}

% -------------------------------------------------------------------------
% Deep Learning / Edge AI Box
% -------------------------------------------------------------------------
\begin{aibox}{สถาปัตยกรรม Edge AI Digital Twin สำหรับการพยากรณ์ความผิดพร่องและการป้องกันอาร์ก}
ปัญหาสำคัญที่สุดในการทำงานของแหล่งจ่ายไฟฟ้าพลาสมาความดันบรรยากาศคือ ปรากฏการณ์เปลี่ยนผ่านสู่การอาร์กอย่างฉับพลัน (Glow-to-Arc Transition) ซึ่งเกิดขึ้นเมื่อความหนาแน่นกระแสเพิ่มขึ้นเฉพาะจุด หากไม่ตัดวงจรภายในไม่กี่ไมโครวินาที พลาสมาอาร์กจะทำลายขั้วไฟฟ้า ไดอิเล็กทริก และชิ้นส่วนพืชอย่างถาวร

\textbf{สถาปัตยกรรมแบบจำลองคู่แฝดดิจิทัล (Digital Twin Architecture):}
คณะผู้วิจัยได้พัฒนาโมเดล Deep 1D-Convolutional Autoencoder (1D-ConvAE) ฝังตัวบนชิปไมโครคอนโทรลเลอร์ Dual-Core สถาปัตยกรรม ESP32-S3 หรือ ARM Cortex-M7 โดยรับสัญญาณกระแสความเร็วสูงจากหัววัด Rogowski Coil ด้วยอัตราสุ่ม $50 \text{ MS/s}$

\begin{equation}
    \mathcal{L}_{\text{AE}}(\mathbf{x}) = \|\mathbf{x} - \hat{\mathbf{x}}\|^2 = \sum_{i=1}^M \left( x_i - \text{Decoder}(\text{Encoder}(x_i)) \right)^2
\end{equation}

\textbf{กลไกการตรวจจับและการสั่งตัดวงจรในระดับไมโครวินาที (Edge In-situ Interruption):}
\begin{itemize}
    \item โมเดลถูกฝึกสอนด้วยรูปคลื่นกระแสในสภาวะปกติ (Glow Discharge และ Streamer Discharge)
    \item เมื่อเริ่มมีสัญญาณบ่งชี้ความไม่เสถียร (Pre-arc Instabilities) เช่น สัญญาณฮาร์มอนิกความถี่สูงผิดปกติ ค่าความคลาดเคลื่อนในการสร้างภาพใหม่ (Reconstruction Error, $\mathcal{L}_{\text{AE}}$) จะพุ่งสูงเกินเกณฑ์เกาส์เซียน ($> \mu + 3\sigma$)
    \item ชิปประมวลผลจะส่งสัญญาณฮาร์ดแวร์อินเทอร์รัพท์ (Hardware Fast Interrupt) ไปปิดเกตของสวิตช์ SiC ภายในเวลาเพียง $1.8 \ \mu\text{s}$ ซึ่งเร็วกว่าระบบฟิวส์และเบรกเกอร์แบบกลไกดั้งเดิมกว่า 1,000 เท่า ช่วยปกป้องเครื่องปฏิกรณ์ได้อย่างสมบูรณ์แบบ
\end{itemize}
\end{aibox}

% -------------------------------------------------------------------------
% Empirical Agricultural Case Study Box (RBRU Standards)
% -------------------------------------------------------------------------
\begin{agribox}{ผลการทดสอบตู้ควบคุมพลาสมามาตรฐาน IP65 ในสภาพแปลงเกษตรกรรมจริง จ.จันทบุรี}
คณะผู้วิจัย มหาวิทยาลัยราชภัฏรำไพพรรณี ได้ดำเนินการสร้างและนำตู้ควบคุมระบบพลาสมาอัจฉริยะมาตรฐานอุตสาหกรรม (Industrial IP65 Smart Plasma Controller) ไปติดตั้งใช้งานจริงในสภาพแวดล้อมที่มีความชื้นสัมพัทธ์สูง ($> 85\% \text{ RH}$) ณ โรงคัดแยกผลไม้และโรงเพาะชำทุเรียน อำเภอท่าใหม่ จังหวัดจันทบุรี ต่อเนื่องเป็นเวลา 180 วัน โดยทำการทดสอบความปลอดภัยทางวิศวกรรมไฟฟ้าและความเข้ากันได้ทางแม่เหล็กไฟฟ้า

\begin{table}[H]
\centering
\small
\caption{ผลการตรวจสอบมาตรฐานความปลอดภัยทางวิวิศวกรรมไฟฟ้าและการทำงานภาคสนาม ($n = 30$ ครั้งทดสอบ)}
\label{tab:ch10_electrical_safety}
\begin{tabular}{lcccc}
\toprule
\textbf{การทดสอบมาตรฐาน} & \textbf{เกณฑ์มาตรฐานสากล} & \textbf{ผลการทดสอบจริง} & \textbf{สถานะการประเมิน} & \textbf{$p$-value} \\
\midrule
ความทนทานต่อแรงดันไดอิเล็กทริก & $2.5 \text{ kV}_{\text{AC}}$ นาน 60 s & $\mathbf{4.0 \text{ kV}_{\text{AC}}}$ ผ่านสมบูรณ์ & ผ่านเกณฑ์ มอก./IEC & --- \\
ความต้านทานฉนวนไฟฟ้า (Insulation) & $> 100 \ \text{M}\Omega$ (ที่ $1000\text{V}$) & $\mathbf{1,450 \pm 85 \ \text{M}\Omega}$ & ผ่านเกณฑ์ขั้นสูง & $<0.001$ \\
กระแสไฟฟ้ารั่วไหล (Leakage Current) & $< 0.75 \text{ mA}$ & $\mathbf{0.042 \pm 0.005 \text{ mA}}$ & ปลอดภัยยิ่งยวด & $<0.001$ \\
ความต้านทานหลักดิน (Earth Resistance) & $< 5.0 \ \Omega$ & $\mathbf{1.85 \pm 0.12 \ \Omega}$ & ผ่านมาตรฐาน IEEE 80 & $<0.001$ \\
ระดับสัญญาณรบกวน EMI ขอบเขต B & ตามเกณฑ์ CISPR 11 & ต่ำกว่าเส้นขีดจำกัด $12 \text{ dB}$ & ผ่านเกณฑ์ระดับอุตสาหกรรม & --- \\
เวลาตอบสนองการตัดอาร์กของ AI & $< 10.0 \ \mu\text{s}$ & $\mathbf{1.82 \pm 0.15 \ \mu\text{s}}$ & แม่นยำ $100\%$ ($n=50$ ครั้ง) & $<0.001$ \\
โอโซนรั่วไหลในบรรยากาศทำงาน & $< 0.10 \text{ ppm}$ (OSHA) & $\mathbf{0.012 \pm 0.003 \text{ ppm}}$ & ปลอดภัยต่อคนงาน & $<0.001$ \\
\bottomrule
\end{tabular}
\end{table}

\textbf{การตรวจสอบความใช้ได้ของวิธีและการประเมินทางสถิติ:}
\begin{itemize}
    \item \textbf{การทดสอบ Paired Samples $t$-test:} ค่าความปลอดภัยทางไฟฟ้าและระดับโอโซนที่ควบคุมได้ต่ำกว่าเกณฑ์มาตรฐานอย่างมีนัยสำคัญทางสถิติยิ่ง ($t_{\text{stat}} = 42.15, p < 0.001$)
    \item \textbf{การตรวจสอบความเที่ยง (\%RSD)} จากการวัดซ้ำ 11 ครั้ง ($n=11$) ของกระแสไฟฟ้ารั่วไหลและเวลาตอบสนอง มีค่า $\%RSD = 1.94\%$ สอดคล้องตามเกณฑ์ EURACHEM ($\%RSD < 5\%$)
    \item \textbf{ความคงทนต่อสภาพแวดล้อมเขตร้อนชื้น:} ตู้ควบคุมทำงานต่อเนื่อง 180 วันในสภาพความชื้นสูงโดยไม่มีการหยุดชะงักของระบบ ไม่มีสนิมกัดกร่อน และไม่มีการขัดข้องจากการรบกวนของคลื่นแม่เหล็กไฟฟ้าต่อเซนเซอร์และระบบ Wi-Fi/IoT ในบริเวณใกล้เคียง
\end{itemize}
\end{agribox}

\section{สรุปสาระสำคัญประจำบท}

การวิศวกรรมระบบพลาสมาความดันบรรยากาศจำเป็นต้องบูรณาการความรู้ด้านอิเล็กทรอนิกส์กำลังขั้นสูง วิศวกรรมไฟฟ้าแรงสูง และความปลอดภัยในการปฏิบัติงานเข้าด้วยกันอย่างเหนียวแน่น แหล่งจ่ายไฟฟ้าพัลส์นาโนวินาทีที่ใช้สวิตช์สารกึ่งตัวนำ SiC/GaN ช่วยให้สามารถควบคุมพลังงานที่ป้อนเข้าสู่อิเล็กตรอนได้อย่างแม่นยำโดยปราศจากความร้อนสูญเสีย การออกแบบการแมตชิ่งอิมพีแดนซ์และการป้องกันสัญญาณรบกวน EMI ตามมาตรฐาน IEC 61000 มีความสำคัญยิ่งต่อความเสถียรของระบบในสภาพแปลงเกษตรกรรม ขณะที่การประยุกต์ใช้โมเดล Edge AI Digital Twin สำหรับการตรวจจับความผิดปกติช่วยให้สามารถตัดวงจรก่อนเกิดการอาร์กได้ภายใน $1.8 \ \mu\text{s}$ สร้างความมั่นใจในความปลอดภัยต่อทั้งเครื่องจักร ผลิตผลเกษตร และผู้ปฏิบัติงานในระดับสูงสุด

\section{คำถามทบทวนท้ายบทที่ 10}

\begin{enumerate}
    \item จงอธิบายหลักการทำงานของมาร์กซ์เจเนอเรเตอร์สารกึ่งตัวนำ (Solid-State Marx Generator) ในสถานะการประจุแบบขนานและคายประจุแบบอนุกรม
    \item เหตุใดสวิตช์สารกึ่งตัวนำประเภท Silicon Carbide (SiC) และ Gallium Nitride (GaN) จึงเหมาะสมกับการสร้างแหล่งจ่ายไฟฟ้าพลาสมาพัลส์มากกว่าสวิตช์ซิลิคอนดั้งเดิม
    \item จงอธิบายปรากฏการณ์การเปลี่ยนแปลงอิมพีแดนซ์ของโหลด DBD ระหว่างช่วงก่อนเบรกดาวน์กับช่วงหลังเบรกดาวน์ และแนวทางการออกแบบเครือข่ายแมตชิ่งเพื่อลดทอนกำลังงานสะท้อนกลับ
    \item แหล่งจ่ายไฟฟ้าพัลส์มาร์กซ์เจเนอเรเตอร์จำนวน 10 สเตจ ใช้ตัวเก็บประจุประจำสเตจ $C_{\text{stage}} = 10 \text{ nF}$ ประจุด้วยแรงดัน $V_{\text{DC}} = 1,200 \text{ V}$ จงคำนวณหาแรงดันพัลส์รวม พลังงานสะสมต่อพัลส์ และกำลังไฟฟ้าเฉลี่ยที่ความถี่ 20 kHz
    \item จงอธิบายสถาปัตยกรรมของ Edge AI Autoencoder ในการตรวจจับปรากฏการณ์ Glow-to-Arc Transition และกระบวนการป้องกันความเสียหายในระดับไมโครวินาที
\end{enumerate}
